\newpage
\section{Simulations réalisées -- 3D}

Cette section présente les simulations numériques directes réalisées en trois dimensions,
d'abord sur une bulle isolée en ascension libre, puis sur une bulle évoluant au sein d'un
écoulement turbulent forcé.

% ─────────────────────────────────────────────────────────────────────────────
\subsection{La bulle isolée (\textit{single rising bubble}) -- 3D}

Des tests préliminaires ont été exécutés en local. La Figure~\ref{fig:snaps3d} illustre
l'évolution temporelle de l'interface. Une rupture de la bulle est observée à la quatrième
seconde de simulation.

\begin{figure}[htbp]
    \centering
    \begin{subfigure}{0.19\textwidth}
        \includegraphics[width=\textwidth]{bulle3d_00.png}
        \caption{$t=0$}
    \end{subfigure}
    \hfill
    \begin{subfigure}{0.19\textwidth}
        \includegraphics[width=\textwidth]{bulle3d_01.png}
        \caption{$t=1$}
    \end{subfigure}
    \hfill
    \begin{subfigure}{0.19\textwidth}
        \includegraphics[width=\textwidth]{bulle3d_02.png}
        \caption{$t=2$}
    \end{subfigure}
    \hfill
    \begin{subfigure}{0.19\textwidth}
        \includegraphics[width=\textwidth]{bulle3d_03.png}
        \caption{$t=3$}
    \end{subfigure}
    \hfill
    \begin{subfigure}{0.19\textwidth}
        \includegraphics[width=\textwidth]{bulle3d_04_fracture.png}
        \caption{$t=4$ (Rupture)}
    \end{subfigure}
    \caption{Évolution temporelle de l'interface en 3D. Maillage grossier (\texttt{maxlevel = 7}).}
    \label{fig:snaps3d}
\end{figure}

Cette fragmentation prématurée résulte de la combinaison de contraintes physiques et numériques.
D'une part, le nombre de Bond est fixé à $Bo = 13,6$, plaçant la dynamique dans le régime des
calottes sphériques avec jupe l'arrière de la bulle s'aplatit sous l'effet de la poussée
d'Archimède. D'autre part, la résolution spatiale est limitée à \texttt{maxlevel = 7}, soit
environ 32 mailles sur le diamètre, insuffisant pour stabiliser l'interface dans ce régime.
Un raffinement à \texttt{maxlevel = 8} ou \texttt{9} est nécessaire, ce qui motive le recours
au supercalculateur pour la suite.

% ─────────────────────────────────────────────────────────────────────────────
\subsection{Bulle en milieu turbulent}

\subsubsection{Stratégie : une DNS résolue par raffinement adaptatif}

L'objectif est de quantifier l'effet de la turbulence ambiante sur l'ascension de la bulle.
Contrairement aux écoulements industriels à très grand Reynolds (où une DNS coûterait
$\sim Re^{9/4}$ mailles et devient prohibitive), le régime visé ici est \emph{modéré}
($Re_\lambda \approx 45$). Une \textbf{simulation numérique directe} qui résout \emph{toutes}
les échelles, de l'échelle intégrale jusqu'à l'échelle dissipative de Kolmogorov $\eta$,
y est donc accessible. La difficulté n'est pas le nombre de Reynolds mais le coût du
\emph{diphasique} (fort ratio de densité $\rho_L/\rho_G = 850$, tension de surface). La
stratégie retenue est une DNS dont la résolution est \textbf{concentrée par raffinement
adaptatif} (\texttt{adapt\_wavelet}) là où les petites échelles vivent, c'est-à-dire sur
la vorticité et sur l'interface. On montre plus bas que le niveau maximal
\texttt{maxlevel = 8} suffit à résoudre $\eta$.

\subsubsection{Protocole précurseur $\rightarrow$ injection}

Le solveur (\texttt{src/main.c}) fonctionne en deux modes issus d'un même code, sélectionnés
par l'argument \texttt{INJECT} :

\begin{enumerate}
    \item \textbf{Précurseur (\texttt{INJECT=0})} -- monophasique ($f=1$ partout). Le champ
    est initialisé par une superposition de modes ABC ($k=1,2,3$) plus un bruit de $10\,\%$,
    puis développe une turbulence homogène isotrope (HIT) \emph{forcée} jusqu'à un état
    statistiquement stationnaire. Le forçage est un \emph{forçage linéaire contrôlé}
    (Bassenne \textit{et al.}, 2016) : à chaque pas, l'amplitude est ajustée pour que la
    puissance injectée compense exactement la dissipation, ce qui \textbf{fixe l'énergie
    cinétique} à une cible \texttt{KE\_TARGET} choisie. L'intensité de la turbulence est
    ainsi un paramètre \emph{piloté} (c'est la grandeur que l'étude fait varier). Le champ
    mûr est sauvegardé (\texttt{end}).
    \item \textbf{Bulle (\texttt{INJECT=1})} -- le champ turbulent précédent est restauré
    (\texttt{restore}), une coquille est raffinée autour de la position de la future interface
    \emph{avant} l'appel à \texttt{fraction()}, puis la sphère VOF est injectée dans une
    turbulence \emph{déjà mature} et laissée remonter.
\end{enumerate}

L'intérêt est d'injecter la bulle dans une turbulence réaliste et stationnaire, et non dans
un transitoire artificiel. Une remarque physique importante : une \emph{unique} mode ABC est
une solution de Beltrami stationnaire et ne transitionne jamais vers la turbulence ; c'est la
superposition multi-modes incommensurables qui déclenche la cascade non linéaire.

\subsubsection{Évolution du maillage : la première approche était défaillante}

La première série de calculs (Tableaux~\ref{tab:convergence_bulle} et
\ref{tab:turbu_bulle_compare}) présentait plusieurs défauts \textbf{rédhibitoires} qui ont
invalidé l'étude de convergence et motivé une réécriture du solveur :

\begin{itemize}
    \item \textbf{Critère d'adaptation trop laxiste.} Le seuil d'ondelette ne forçait jamais
    le raffinement de l'interface au niveau maximal : les niveaux 7, 8 et 9 produisaient le
    \emph{même nombre effectif de mailles}, si bien que la « convergence » ne testait rien.
    \item \textbf{Perte de volume à l'injection.} Faute de mailles fines à l'emplacement de la
    bulle lors de \texttt{fraction()}, la bulle perdait plus de $50\,\%$ de son volume
    théorique dès l'injection (colonne $V/V_{th}$ du Tableau~\ref{tab:convergence_bulle} :
    $1\,\%$, $5\,\%$, $45\,\%$).
    \item \textbf{Rayon trop petit} ($R_0 = 4{,}24$) : $\sim 17$ mailles par diamètre au
    niveau 7, sous le seuil admis de $20$.
    \item \textbf{Précurseur non stationnaire} au niveau 9 (spin-up trop court,
    $\mathrm{d}k/\mathrm{d}t < 0$), rendant les grandeurs turbulentes non comparables entre
    niveaux, avec de surcroît des instants d'injection désynchronisés.
\end{itemize}

Ces tableaux sont conservés ci-dessous à titre de référence historique.

\begin{table}[H]
    \centering
    \caption{[\textbf{Première approche, défaillante}] Convergence spatiale. Volume théorique
    $V_{th} = \frac{4}{3}\pi R_0^3 \approx 333$ pour $R_0 = 4.3$. Grandeurs turbulentes sur le
    précurseur ($t \in [120,140]$).}
    \label{tab:convergence_bulle}
    \begin{tabular}{lccc}
        \toprule
        \textbf{Grandeur} & \textbf{lvl 7} & \textbf{lvl 8} & \textbf{lvl 9} \\
        \midrule
        $\Delta x$ [-]               & $0.9375$ & $0.4688$ & $0.2344$ \\
        \midrule
        $k$ [-]                      & $\approx 48.7$ & $\approx 49.3$ & $\approx 36.5^{\star}$ \\
        $\varepsilon$ [-]            & $\approx 0.46$ & $\approx 0.47$ & $\approx 0.46^{\star}$ \\
        $Re_\lambda$ [-]             & $\approx 137$  & $\approx 137$  & $\approx 102^{\star}$ \\
        \midrule
        Volume bulle [-]             & $\approx 3.3$  & $\approx 17$   & $\approx 150$ \\
        $V/V_{th}$ [\%]              & $\approx 1.0$  & $\approx 5.1$  & $\approx 45$ \\
        Nb fragments                 & $1$ & $1$ & $1$ \\
        \bottomrule
    \end{tabular}

    \vspace{0.4em}
    {\footnotesize $^{\star}$ Précurseur non stationnaire à lvl 9 (spin-up trop court,
    $\mathrm{d}k/\mathrm{d}t < 0$ à l'injection) : valeurs non comparables.}
\end{table}

\begin{table}[H]
    \centering
    \caption{[\textbf{Première approche}] Grandeurs turbulentes du liquide avant et après
    injection de la bulle (lvl 8).}
    \label{tab:turbu_bulle_compare}
    \begin{tabular}{lcc}
        \toprule
        \textbf{Grandeur} & \textbf{Sans bulle} ($t \in [120,199]$) & \textbf{Avec bulle} ($t \in [200,250]$) \\
        \midrule
        $k$ [-]            & $\approx 49.3$  & $\approx 49.35$ \\
        $\varepsilon$ [-]  & $\approx 0.469$ & $\approx 0.469$ \\
        $Re_\lambda$ [-]   & $\approx 137.4$ & $\approx 137.8$ \\
        \bottomrule
    \end{tabular}
\end{table}

\subsubsection{Maillage corrigé et validé}

La réécriture apporte quatre corrections, résumées dans le Tableau~\ref{tab:corrections_maillage} :

\begin{table}[H]
    \centering
    \caption{Corrections apportées au maillage et à l'initialisation.}
    \label{tab:corrections_maillage}
    \begin{tabular}{@{}p{4.2cm}p{4.2cm}p{6.3cm}@{}}
    \toprule
    \textbf{Défaut (v1)} & \textbf{Correction (v2)} & \textbf{Effet} \\ \midrule
    Ondelette trop lâche & Seuil interface $f_{\max}=10^{-3}$ & Force l'interface au
    \texttt{maxlevel} : les niveaux deviennent réellement distincts. \\
    Perte de volume ($>50\,\%$) & Coquille raffinée avant \texttt{fraction()} & Volume
    conservé à $97$--$99\,\%$ à l'injection (v. ci-dessous). \\
    $R_0 = 4{,}24$ ($<20$ mailles/$D$) & $R_0 = 8$ & $D/\Delta x \approx 34$ au niveau 8. \\
    Adaptation sur la vitesse & Adaptation sur $\{f,\,|\boldsymbol\omega|\}$ & Raffine les
    petites échelles (vorticité) et l'interface, pas les grosses structures inutiles. \\
    \bottomrule
    \end{tabular}
\end{table}

Pour le précurseur, la grille est \textbf{uniforme} pendant le développement
($t < 18$, pour ne pas biaiser l'adaptation par le bruit transitoire), puis passe en
\textbf{AMR sur la vorticité}. Cette AMR répond aussi à une contrainte machine : sur la
version Basilisk du calculateur, \texttt{restore()} d'un cube $256^3$ uniforme
($16{,}8$~M mailles) provoque un défaut de segmentation. L'AMR ramène le champ à
$\mathbf{6{,}6}$~\textbf{M mailles} (dump $550$~Mo au lieu de $1{,}45$~Go), qui se restaure
correctement, tout en \emph{améliorant} la résolution de la turbulence par rapport à un
$128^3$. La conservation de volume à l'injection, désormais de $94$, $98$ et $99\,\%$ pour
les niveaux $7$, $8$ et $9$, confirme la correction (elle \emph{croît} avec la résolution,
comme attendu).

\subsubsection{Caractérisation de la turbulence homogène isotrope}
\label{sec:carac_hit_ref}

La Figure~\ref{fig:prec_carac} montre le développement du champ. La dissipation $\varepsilon$
part de la valeur d'un champ lisse ($\approx 6$), \textbf{croît} à mesure que la cascade
non linéaire s'installe (signature de la transition), culmine vers $t\approx 18$, puis se
stabilise. $Re_\lambda$, $k$ et $\eta$ atteignent un plateau : l'état est statistiquement
stationnaire pour $t \gtrsim 22$ (zone verte), fenêtre sur laquelle les moyennes sont prises.

\begin{figure}[H]
    \centering
    \includegraphics[width=0.92\textwidth]{prec_carac.png}
    \caption{Développement de la HIT du précurseur (level-8-AMR). Énergie $k$, dissipation
    $\varepsilon$, Reynolds de Taylor $Re_\lambda$ et échelle de Kolmogorov $\eta$ en fonction
    du temps. Trait pointillé : début de l'AMR ($t=18$) ; bande verte : fenêtre stationnaire.}
    \label{fig:prec_carac}
\end{figure}

\begin{table}[H]
    \centering
    \caption{Caractérisation de la turbulence (moyennes sur $t \ge 22$, unités code :
    $\rho_1=1$, $g=1$, $L_0=120$, $D=2R_0=16$).}
    \label{tab:carac_hit}
    \begin{tabular}{@{}llc@{}}
    \toprule
    \textbf{Grandeur} & \textbf{Symbole} & \textbf{Valeur} \\ \midrule
    Reynolds de Taylor          & $Re_\lambda = u'\lambda/\nu$      & $\mathbf{44{,}8}$ \\
    Reynolds intégral           & $Re_L = u'L/\nu$                  & $246$ \\
    Fluctuation rms             & $u' = \sqrt{2k/3}$                & $10{,}65$ \\
    Énergie cinétique           & $k$                               & $170{,}2$ \\
    Dissipation                 & $\varepsilon$                     & $52{,}8$ \\
    Viscosité cinématique       & $\nu$                             & $1{,}824$ \\ \midrule
    Échelle intégrale           & $L = k^{3/2}/\varepsilon$         & $42{,}1$ \\
    Micro-échelle de Taylor     & $\lambda$                         & $7{,}67$ \\
    Échelle de Kolmogorov       & $\eta = (\nu^3/\varepsilon)^{1/4}$& $0{,}582$ \\
    Temps de retournement       & $\tau_e = k/\varepsilon$          & $3{,}23$ \\
    Temps de Kolmogorov         & $\tau_\eta = (\nu/\varepsilon)^{1/2}$ & $0{,}186$ \\ \midrule
    Maille fine (niveau 8)      & $\Delta x = L_0/2^8$              & $0{,}469$ \\
    \textbf{Résolution}         & $k_{\max}\,\eta$                  & $\mathbf{3{,}90}$ \\
    Kolmogorov / maille         & $\eta/\Delta x$                   & $1{,}24$ \\
    \bottomrule
    \end{tabular}
\end{table}

\paragraph{Justification de la résolution -- toutes les échelles sont résolues.}
Le critère standard de résolution d'une DNS (Pope) est $k_{\max}\,\eta \gtrsim 1{,}5$, avec
$k_{\max} = \pi\, 2^{8}/L_0$ le plus grand nombre d'onde résolu au niveau 8. On obtient
$k_{\max}\,\eta = \mathbf{3{,}90}$, très au-dessus du seuil : l'échelle de Kolmogorov est
non seulement atteinte mais \emph{sur-résolue} ($\eta \approx 1{,}24\,\Delta x$, soit plus
d'une maille par échelle dissipative). La gamme dissipative -- donc l'intégralité du spectre
d'énergie -- est capturée. La visualisation du maillage (Figure~\ref{fig:prec_visu}, coin
supérieur) confirme que l'AMR place effectivement les mailles fines sur les filaments de
vorticité et laisse grossir les zones lisses, ce qui rend la résolution $\eta$ atteignable à
coût réduit.

\paragraph{Résolution de la couche limite de la bulle.}
Au voisinage de l'interface, le nombre de Reynolds de bulle vaut
$Re_b = U_b D/\nu \approx 35$ (avec $U_b = \sqrt{gD} = 4$). L'épaisseur de couche limite
attendue $\delta \sim D/\sqrt{Re_b} \approx 2{,}7$ est donc résolue par
$\delta/\Delta x \approx 6$ mailles au niveau 8 (davantage au niveau 9), ce qui est suffisant
pour capturer le cisaillement pariétal et le sillage proche.

\paragraph{Changement de nombre de Reynolds et coût.}
Le $Re_\lambda \approx 45$ actuel est \emph{plus faible} que la valeur $\approx 137$ de la
première approche (par ailleurs non fiable car non stationnaire). Ce choix est \textbf{dicté
par la résolution} : à \texttt{maxlevel} fixé, résoudre $\eta$ impose une borne supérieure
au Reynolds atteignable (un $Re_\lambda$ plus grand $\Rightarrow$ $\eta$ plus petit
$\Rightarrow$ maillage plus fin). À $Re_\lambda = 45$, $\eta = 0{,}582$ tombe pile dans les
capacités du niveau 8 ($k_{\max}\eta = 3{,}9$) ; viser $Re_\lambda \sim 70$--$80$ exigerait
un niveau 9--10, incompatible avec la contrainte \texttt{restore} et le coût. Le
\texttt{maxlevel = 8} avec AMR-vorticité est donc le point d'équilibre optimal : DNS
pleinement résolue au coût le plus bas. Le développement du précurseur ($t=0 \to 30$,
$5590$ pas) a coûté $\approx 5{,}4$~h sur $64$ cœurs ($\sim 345$~h.CPU), avec $6{,}6$~M mailles
adaptatives au lieu de $16{,}8$~M en uniforme ($\times 2{,}5$ d'économie).

\paragraph{Régime de la bulle.} L'intensité turbulente relative à la bulle est élevée,
$I = u'/U_b = 2{,}66$ : les fluctuations turbulentes dominent la flottabilité. Le nombre de
Weber turbulent, $We_t = 2\rho_1 (\varepsilon D)^{2/3} D/\sigma \approx 11$, place la bulle
dans un régime de \emph{forte déformation} (voire de fragmentation), cohérent avec les
rapports d'aspect $\chi \gtrsim 2$ observés. Ce point sera quantifié sur les runs bulles en cours.

\subsubsection{Visualisation des structures turbulentes}

La Figure~\ref{fig:prec_visu} présente le champ mûr ($t=25$) : le maillage adaptatif suivant
les structures, la vorticité $|\boldsymbol\omega|$ et la vitesse verticale $u_z$ sur une coupe
à mi-domaine, ainsi qu'une isosurface 3D de $|\boldsymbol\omega|$ (seuil $1{,}5\,\omega_{rms}$)
colorée par $u_z$, révélant les tubes tourbillonnaires en « vermicelles » caractéristiques de
la HIT. Des films de l'évolution temporelle ($t \in [19,30]$) du champ de vitesse, de la
vorticité et de l'isosurface 3D ont été produits (\texttt{scripts/figures/films/}).

\begin{figure}[H]
    \centering
    \begin{subfigure}[b]{0.48\textwidth}
        \includegraphics[width=\textwidth]{prec_maillage.png}
        \caption{Maillage adaptatif (coupe $z$)}
    \end{subfigure}
    \hfill
    \begin{subfigure}[b]{0.48\textwidth}
        \includegraphics[width=\textwidth]{prec_vortex3d.png}
        \caption{Isosurface 3D de $|\boldsymbol\omega|$, colorée par $u_z$}
    \end{subfigure}

    \vspace{0.5em}
    \begin{subfigure}[b]{0.48\textwidth}
        \includegraphics[width=\textwidth]{prec_vorticite.png}
        \caption{Vorticité $|\boldsymbol\omega|$ (coupe $z$)}
    \end{subfigure}
    \hfill
    \begin{subfigure}[b]{0.48\textwidth}
        \includegraphics[width=\textwidth]{prec_vitesse.png}
        \caption{Vitesse verticale $u_z$ (coupe $z$)}
    \end{subfigure}
    \caption{Précurseur turbulent mûr (level-8-AMR, $t=25$, $Re_\lambda \approx 45$).
    Le maillage se resserre sur les filaments de forte vorticité ; les structures
    tourbillonnaires 3D sont lisses, donc bien résolues.}
    \label{fig:prec_visu}
\end{figure}

\subsubsection{Premiers résultats bulle : rupture turbulente observée (lvl 7, \texttt{KE\_TARGET}=200)}
\label{sec:resultats_lvl7}

Le run \texttt{lvl7} -- bulle injectée dans le précurseur ci-dessus à $t=30$ et suivie
jusqu'à $t=80$, soit $50$ unités de temps physique -- vient de se terminer sur Kairos
($32$ cœurs, $4$h$46$ de calcul). Le solveur lui-même n'a montré aucun signe de
faiblesse : nombre de mailles stable ($1{,}36$--$1{,}70$~M), pas de temps stable
($7{,}3\times10^{-4}$--$1{,}3\times10^{-2}$), itérations de multigrille pression
plafonnant à $45$, très en dessous du seuil de blocage ($\approx 100$) rencontré par
ailleurs au niveau 8. À l'injection, le volume vaut $2021$, soit $94\,\%$ du volume
théorique $\frac{4}{3}\pi R_0^3 \approx 2145$ -- cohérent avec la conservation
annoncée au Tableau~\ref{tab:corrections_maillage}.

En revanche, la bulle \textbf{ne reste pas intacte}. Le champ turbulent d'injection
est celui caractérisé plus haut ($We_t \approx 11$, régime annoncé de forte
déformation « voire de fragmentation », \S\ref{sec:carac_hit_ref}) : ce point,
laissé en suspens, est maintenant quantifié par le Tableau~\ref{tab:rupture_lvl7}.

\begin{table}[H]
    \centering
    \caption{Chronologie de la rupture, reconstruite depuis \texttt{bubble.dat}. Le
    « volume » est celui de la plus grande région gazeuse taguée
    (\texttt{tag.h}) -- la « bulle mère » -- pas le volume total de gaz.}
    \label{tab:rupture_lvl7}
    \begin{tabular}{@{}lp{9.7cm}@{}}
        \toprule
        \textbf{Instant} & \textbf{Observation} \\
        \midrule
        $t=30{,}0$ (injection) & Volume $=2021$ ($94\,\%$ du théorique) ; $d_{eq}=15{,}7 \approx 2R_0$. \\
        $t=34{,}1$ & Étirement maximal : la boîte englobante atteint $44{,}1$ sur un axe,
        $\approx 3\times$ le diamètre injecté -- signature d'un ligament juste avant
        pincement. \\
        $t\approx 38$--$41$ & \textbf{Rupture.} Le volume de la région mère chute de
        $\approx 1335$ (moyenne $t<40$) à $\approx 250$ (moyenne $t>45$), soit
        $\approx 88\,\%$ de perte apparente ; $d_{eq}$ final $=7{,}3$, la moitié du
        diamètre injecté. \\
        $t=46{,}1$ & Pic de déformation : $\chi = 4{,}60$ (rapport d'aspect max/min) --
        un fragment encore en cours de déchirement. \\
        $t=30 \to 80$ & Sauts de volume ponctuels de $\pm20$--$39\,\%$ en un seul pas
        d'écriture (ex. $t=65{,}6$--$65{,}9$ : $211\to265\to211\to208$) : le tag
        « mère » bascule entre fragments de taille comparable -- bruit de
        ré-étiquetage plutôt qu'évolution physique continue. \\
        \bottomrule
    \end{tabular}
\end{table}

\begin{figure}[H]
    \centering
    \includegraphics[width=0.75\textwidth]{bulle_lvl7_uz.png}
    \caption{Vitesse de remontée $u_z(t)$ de la région mère, lvl 7. Aucun plateau :
    $u_z$ oscille entre $-17{,}7$ et $+34{,}3$ (écart-type $9{,}2$ autour d'une
    moyenne $0{,}9$) sur toute la fenêtre de suivi.}
    \label{fig:lvl7_uz}
\end{figure}

\begin{figure}[H]
    \centering
    \includegraphics[width=0.95\textwidth]{bulle_lvl7_aspect.png}
    \caption{Rapport d'aspect $\chi(t)$ et dimensions de la boîte englobante, lvl 7.
    La chute des trois dimensions après $t\approx 40$ ($44 \to 5$--$10$) marque la
    rupture.}
    \label{fig:lvl7_aspect}
\end{figure}

\begin{figure}[H]
    \centering
    \includegraphics[width=0.95\textwidth]{bulle_lvl7_trajectoire.png}
    \caption{Montée $z(t)-z(0)$ (gauche) et dérive horizontale $x$--$y$ (droite),
    lvl 7. La montée nette est positive ($+87$ sur $50$ unités) mais passe par un
    pic à $+162$ ($t=57{,}5$) suivi d'une redescente : pas de vitesse terminale au
    sens classique. Les segments rectilignes traversant tout le domaine périodique
    à droite sont la signature d'un changement d'identité de fragment, non d'une
    dérive physique continue.}
    \label{fig:lvl7_trajectoire}
\end{figure}

\paragraph{Limite méthodologique du suivi « bulle mère ».} Le commentaire du code
(\texttt{event bubble}, \texttt{src/main.c}) suppose explicitement qu'à $Bo=1$ « pas
de fragmentation n'est attendue », et ne trace donc que la région gazeuse la plus
volumineuse. Cette hypothèse ne tient plus à $We_t \approx 11$ : après rupture, la
« vitesse terminale » que produirait une moyenne temporelle sur la fin du signal
(ex. $-1{,}95$ sur les $20\,\%$ derniers instants) ne décrit qu'un fragment parmi
d'autres, et non la bulle physique injectée -- elle n'a donc pas de sens comme point
de comparaison de traînée. Une amélioration possible, si la rupture réapparaît sur
les niveaux suivants, serait de tracer en parallèle le volume gazeux \emph{total}
(somme sur tous les tags) afin de distinguer perte physique réelle et bruit de
ré-étiquetage.

\paragraph{Conséquence sur l'étude de convergence.} Le niveau 8 s'est révélé, à
$KE_{target}=200$, trop coûteux pour converger (le multigrille pression sature son
plafond d'itérations à chaque pas, cf. \S\ref{sec:carac_hit_ref}) : son précurseur a
donc été \textbf{ré-équilibré} vers $KE_{target}=50$ ($We_t \approx 2$--$3$,
Re$_\lambda \approx 32$) par une relaxation monophasique courte avant réinjection de
la bulle. Les runs \texttt{lvl7} ($KE_{target}=200$) et \texttt{lvl8}
($KE_{target}=50$) ne correspondent donc plus à la même physique, et \textbf{ne
peuvent pas être comparés directement} comme un test de convergence spatiale à
intensité fixée. Le run \texttt{lvl8} sur champ ré-équilibré à $KE_{target}=50$ a
cependant été interrompu dès l'injection ($t=48$, avant tout développement physique
de la bulle) au profit d'un second ré-équilibrage, décrit ci-dessous, visant à
éliminer la fragmentation plutôt que de composer avec elle.

\subsubsection{Second ré-équilibrage : $We_t \approx 1$ pour éliminer la fragmentation}
\label{sec:carac_we1}

À $We_t \approx 2$--$3$ (champ $KE_{target}=50$), rien ne garantissait que la bulle
lvl8 resterait intacte : le régime restait, par construction, celui d'une
« déformation modérée » et non celui, visé pour une étude de convergence propre,
d'une bulle qui garde son identité sur toute la fenêtre de suivi. Le champ a donc
été relaxé une seconde fois, \textbf{depuis le \texttt{dump} du ré-équilibrage
précédent} ($t=48$, $KE_{target}=50$), vers une cible plus conservatrice
$KE_{target}=15$ (\texttt{./bubble 8 108 8.0 1 15 0}, toujours monophasique).

\begin{figure}[H]
    \centering
    \includegraphics[width=0.92\textwidth]{prec_carac_we1.png}
    \caption{Second ré-équilibrage monophasique, niveau 8 : décroissance de
    $KE_{target}=50$ vers $KE_{target}=15$ ($t=48 \to 77$). Trait pointillé :
    reprise du dump ($t=48$) ; bande verte : fenêtre de moyenne ($t \ge 67$).}
    \label{fig:prec_carac_we1}
\end{figure}

Contrairement au précurseur initial (démarrage à froid, transition bruit
$\to$ turbulence), il s'agit ici d'une simple relaxation d'un champ déjà turbulent
vers une énergie plus basse : $k$ et $\varepsilon$ décroissent de façon monotone
(Figure~\ref{fig:prec_carac_we1}) et $Re_\lambda$, $\eta$ se stabilisent bien avant
l'instant final. Le run a été \textbf{arrêté à $t=76$} (un instantané entier,
proprement écrit) plutôt que d'attendre $t=108$ comme prévu par le script : $k=14,97$
y est déjà à moins de $2\,\%$ de la cible $15$, prolonger n'aurait fait que consommer
des heures de calcul sans changer la physique (le ré-équilibrage étant monophasique,
le coût reste faible : $\approx 1$h sur $64$ cœurs pour les $29$ unités parcourues).

\begin{table}[H]
    \centering
    \caption{Caractérisation de la turbulence après le second ré-équilibrage
    (moyennes sur $t \in [67, 77]$, niveau 8, mêmes unités code que le
    Tableau~\ref{tab:carac_hit}).}
    \label{tab:carac_we1}
    \begin{tabular}{@{}llc@{}}
    \toprule
    \textbf{Grandeur} & \textbf{Symbole} & \textbf{Valeur} \\ \midrule
    Reynolds de Taylor          & $Re_\lambda = u'\lambda/\nu$      & $\mathbf{24{,}6}$ \\
    Reynolds intégral           & $Re_L = u'L/\nu$                  & $74{,}3$ \\
    Fluctuation rms             & $u' = \sqrt{2k/3}$                & $3{,}16$ \\
    Énergie cinétique           & $k$                               & $14{,}97$ \\
    Dissipation                 & $\varepsilon$                     & $1{,}35$ \\
    Viscosité cinématique       & $\nu$                             & $1{,}824$ \\ \midrule
    Échelle intégrale           & $L = k^{3/2}/\varepsilon$         & $42{,}9$ \\
    Micro-échelle de Taylor     & $\lambda$                         & $14{,}2$ \\
    Échelle de Kolmogorov       & $\eta = (\nu^3/\varepsilon)^{1/4}$& $1{,}456$ \\
    Temps de retournement       & $\tau_e = k/\varepsilon$          & $11{,}08$ \\
    Temps de Kolmogorov         & $\tau_\eta = (\nu/\varepsilon)^{1/2}$ & $1{,}162$ \\ \midrule
    Maille fine (niveau 8)      & $\Delta x = L_0/2^8$              & $0{,}469$ \\
    \textbf{Résolution}         & $k_{\max}\,\eta$                  & $\mathbf{9{,}76}$ \\
    Kolmogorov / maille         & $\eta/\Delta x$                   & $3{,}11$ \\ \midrule
    Intensité relative          & $I = u'/U_b$                      & $0{,}79$ \\
    \textbf{Weber turbulent}    & $We_t = 2\rho_1(\varepsilon D)^{2/3}D/\sigma$ & $\mathbf{0{,}97}$ \\
    \bottomrule
    \end{tabular}
\end{table}

$We_t = 0{,}97$ : la cible $We_t \approx 1$ est atteinte quasiment exactement. La
résolution n'en souffre pas -- au contraire : une turbulence plus faible a une
échelle de Kolmogorov plus grande ($\eta = 1{,}456$ contre $0{,}582$ au champ
$KE=200$), donc \emph{moins} de mailles fines sont nécessaires pour la résoudre
($k_{\max}\eta = 9{,}76$, contre $3{,}90$ précédemment -- très largement sur-résolu).
C'est cohérent avec la Figure~\ref{fig:prec_visu_we1} : le maillage adaptatif est
visiblement plus grossier que pour le champ $KE=200$ (Figure~\ref{fig:prec_visu}),
et donc bien moins coûteux -- $\approx 1{,}3$~M mailles pour le champ $KE=15$ contre
$6{,}6$~M à $KE=200$ ($\times 5$) -- ce qui a permis de relancer les deux niveaux
avec \emph{moins} de cœurs qu'auparavant (voir plus bas).

\begin{figure}[H]
    \centering
    \begin{subfigure}[b]{0.48\textwidth}
        \includegraphics[width=\textwidth]{prec_maillage_we1.png}
        \caption{Maillage adaptatif (coupe $z$)}
    \end{subfigure}
    \hfill
    \begin{subfigure}[b]{0.48\textwidth}
        \includegraphics[width=\textwidth]{prec_vortex3d_we1.png}
        \caption{Isosurface 3D de $|\boldsymbol\omega|$, colorée par $u_z$}
    \end{subfigure}

    \vspace{0.5em}
    \begin{subfigure}[b]{0.48\textwidth}
        \includegraphics[width=\textwidth]{prec_vorticite_we1.png}
        \caption{Vorticité $|\boldsymbol\omega|$ (coupe $z$)}
    \end{subfigure}
    \hfill
    \begin{subfigure}[b]{0.48\textwidth}
        \includegraphics[width=\textwidth]{prec_vitesse_we1.png}
        \caption{Vitesse verticale $u_z$ (coupe $z$)}
    \end{subfigure}
    \caption{Champ précurseur après le second ré-équilibrage ($t=76$,
    $Re_\lambda \approx 24{,}6$, $We_t \approx 1$). À comparer à la
    Figure~\ref{fig:prec_visu} ($KE=200$) : structures tourbillonnaires
    plus grosses et moins nombreuses, maillage nettement plus grossier.}
    \label{fig:prec_visu_we1}
\end{figure}

\paragraph{Relance des bulles lvl7 et lvl8 sur le champ $We_t \approx 1$.} Les runs
\texttt{lvl7}-$KE200$ (fragmenté) et \texttt{lvl8}-$KE50$ (interrompu à l'injection)
sont archivés (\texttt{lvl7\_ke200\_fragmented}, \texttt{lvl8\_ke50\_aborted}) et
remplacés par une injection propre sur ce nouveau champ (\texttt{end\_we1}, $t=76$) :
même protocole (\texttt{INJECT=1}, $R_0=8$, forçage contrôlé maintenu à
$KE_{target}=15$ pendant la remontée), fenêtre de suivi de $50$ unités
($MAXTIME=126$). À l'injection, le volume vaut $93{,}6\,\%$ (lvl7) et $97{,}4\,\%$
(lvl8) du volume théorique, rapport d'aspect $\chi \approx 1$ sur les deux niveaux --
cohérent avec la conservation de volume croissante avec la résolution déjà établie
(Tableau~\ref{tab:corrections_maillage}) et, contrairement au champ $KE=200$, sans
signe de déformation excessive dès l'injection. Le maillage plus grossier de ce
champ (Figure~\ref{fig:prec_visu_we1}) permet de réduire le nombre de cœurs par
rapport aux runs précédents : \texttt{lvl7} tourne sur $8$ cœurs et \texttt{lvl8}
sur $32$ cœurs (walltime $72$h chacun), à comparer aux $32$/$64$ cœurs utilisés pour
les tentatives $KE200$/$KE50$. Ces deux runs sont désormais terminés ; leurs
résultats font l'objet des deux sous-sections suivantes.

\subsubsection{Résultats bulle à $We_t \approx 1$ : bulle intacte et convergence
en maillage}
\label{sec:resultats_we1}

Contrairement au run $KE200$ (\S\ref{sec:resultats_lvl7}), la bulle \textbf{reste
intacte} sur toute la fenêtre de suivi aux deux niveaux. Les deux runs ont été
suivis sur $50$ unités de temps ($t=76 \to 126$) ; le \texttt{lvl8}, d'abord trop
lent, n'a abouti qu'après avoir relâché la tolérance du solveur multigrille à
$0{,}1$ -- accélération sans perte, puisqu'il s'accorde ensuite à $1\,\%$ près avec
le \texttt{lvl7} (\emph{cf.} ci-dessous). La conservation du volume est excellente et \emph{stable}
tout du long : $\approx 94\,\%$ du volume théorique (lvl7) et $\approx 97\,\%$
(lvl8), sans la chute brutale observée à la rupture du run $KE200$
(Tableau~\ref{tab:rupture_lvl7}). Le rapport d'aspect oscille sous
$\chi \lesssim 2{,}5$ sans emballement (contre $\chi = 4{,}6$ au moment de la
rupture $KE200$) : la bulle se \emph{déforme} sous les tourbillons mais ne
fragmente pas -- le régime $We_t \approx 1$ visé est bien celui où la vitesse de
remontée a un sens.

\begin{figure}[H]
    \centering
    \includegraphics[width=\textwidth]{bulle_we1_uz_volume.png}
    \caption{Runs bulle $We_t \approx 1$ (lvl7 et lvl8 \emph{terminés},
    $50$~u.t. chacun). \textbf{Gauche} : vitesse verticale $u_z$ --
    trait fin = instantané (fortement oscillant, la bulle est brassée par les
    tourbillons, écart-type $\approx 2{,}5$--$3$), trait épais = moyenne cumulée
    depuis l'injection. \textbf{Droite} : volume rapporté au volume théorique
    $\frac{4}{3}\pi R_0^3$ -- \textbf{stable} aux deux niveaux, aucune
    fragmentation.}
    \label{fig:bulle_we1_uz}
\end{figure}

\paragraph{Vitesse de remontée : convergence spatiale lvl7 $\approx$ lvl8.}
La vitesse de remontée se lit sur le \emph{déplacement net} (position déroulée du
centre de masse, robuste à la périodicité, cf. \S\ref{sec:carac_we1} et la moyenne
circulaire du code). Sur la fenêtre \textbf{complète} de $50$ unités de temps,
\[
   U_b^{\text{lvl7}} = 1{,}08 \quad\approx\quad U_b^{\text{lvl8}} = 1{,}07
   \qquad (\text{écart } \approx 1\,\%),
\]
et l'accord est déjà de $1\,\%$ sur les $19$ premières unités
($U_b = 2{,}79$ contre $2{,}82$). La vitesse de remontée est donc
\textbf{indépendante du maillage} à intensité turbulente fixée -- exactement le
test de convergence spatiale que la première approche (\S\ref{sec:resultats_lvl7}
et Tableau~\ref{tab:convergence_bulle}) ne pouvait pas fournir. C'est le premier
point de convergence propre de l'étude.

\begin{figure}[H]
    \centering
    \includegraphics[width=\textwidth]{bulle_we1_aspect_traj.png}
    \caption{Runs bulle $We_t \approx 1$. \textbf{Gauche} : rapport d'aspect
    $\chi(t)$ -- oscille sans emballement, $\chi \lesssim 2{,}5$. \textbf{Droite} :
    montée cumulée $z(t)-z_0$ (position déroulée) -- \textbf{superposée} entre lvl7
    et lvl8 sur la fenêtre commune, ce qui confirme visuellement la convergence.}
    \label{fig:bulle_we1_aspect}
\end{figure}

\begin{table}[H]
    \centering
    \caption{Récapitulatif des runs bulle à $We_t \approx 1$
    ($\nu \approx 1{,}82$, $D = d_{eq} \approx 15{,}8$, unités code).}
    \label{tab:bulle_we1}
    \begin{tabular}{@{}lcc@{}}
        \toprule
        \textbf{Grandeur} & \textbf{lvl 7} & \textbf{lvl 8} \\
        \midrule
        Cœurs / fenêtre suivie             & $8$ / $50$~u.t. & $32$ / $50$~u.t. \\
        Cellules (maillage adaptatif)      & $\approx 0{,}66$~M & $\approx 1{,}2$~M \\
        Volume / théorique                 & $\approx 94\,\%$ & $\approx 97\,\%$ \\
        $\chi$ moyen (max)                 & $1{,}26$ ($2{,}44$) & $1{,}24$ ($2{,}50$) \\
        Fragmentation                      & non & non \\ \midrule
        $U_b$ sur $19$~u.t.                & $2{,}79$ & $2{,}82$ \\
        $U_b$ net sur $50$~u.t.            & $\mathbf{1{,}08}$ & $\mathbf{1{,}07}$ \\
        $C_D$ effectif ($Re_b \approx 9$)  & $\approx 18$ & $\approx 18$ \\
        \bottomrule
    \end{tabular}
\end{table}

\paragraph{Nuance : $U_b$ moyen vs vitesse terminale.} La valeur convergée
$U_b \approx 1{,}07$--$1{,}08$ est une \emph{moyenne sur $50$~unités}, dominée par
la montée initiale : la bulle est injectée au repos, accélère fortement pendant les
$\sim 11$ premières unités (pic $u_z \approx 7$, moyenne cumulée $\approx 2{,}8$ sur
les $19$ premières), puis le brassage turbulent prend le dessus et la montée
\emph{nette} devient quasi nulle sur les $25$ dernières unités ($u_z$ instantané
oscille autour de zéro, écart-type $\approx 2{,}5$--$3$). Cette valeur, bien que
\textbf{convergée en maillage} (lvl7 $\equiv$ lvl8), ainsi que le coefficient de
traînée effectif associé $C_D \approx 18$ ($Re_b \approx 9$), restent donc des
grandeurs \emph{moyennées sur la fenêtre} et non une vitesse terminale stationnaire
au sens classique : une seule bulle dans une seule réalisation de turbulence voit sa
montée dominée par les tourbillons qu'elle rencontre. Une vitesse terminale
« propre » exigerait une moyenne d'ensemble (plusieurs bulles ou réalisations) ou
une fenêtre nettement plus longue.

\subsubsection{Effet de la bulle sur la turbulence : un couplage à sens unique}
\label{sec:turbu_avec_bulle}

Une question sous-jacente à toute l'étude est de savoir si la bulle \emph{rétroagit}
sur la turbulence ambiante (couplage à double sens) ou si elle se contente d'être
\emph{transportée} par un champ qu'elle ne modifie pas (couplage à sens unique). On
la tranche en comparant les grandeurs turbulentes du liquide \textbf{avant}
l'injection (champ ré-équilibré, \S\ref{sec:carac_we1}) et \textbf{pendant} le run
bulle, \emph{au même niveau 8} pour ne comparer que l'effet de la bulle et non celui
de la résolution. La Figure~\ref{fig:turbu_avec_bulle} montre que les séries
temporelles $k$, $\varepsilon$, $Re_\lambda$ sont \textbf{continues} à travers
l'injection : aucune rupture, aucun régime transitoire déclenché par l'apparition
de la bulle.

\begin{figure}[H]
    \centering
    \includegraphics[width=\textwidth]{turbu_avec_bulle.png}
    \caption{Grandeurs turbulentes du liquide forcé ($KE_{target}=15$, niveau 8) de
    part et d'autre de l'injection de la bulle ($t=76$, trait pointillé vertical).
    Bleu : champ ré-équilibré \emph{sans} bulle ; rouge : run \texttt{lvl8}
    \emph{avec} bulle ; pointillé vert : moyenne de référence sans bulle. Les
    courbes se raccordent sans saut ; l'écart après injection reste modéré
    ($\lesssim 10\,\%$), avec une lente dérive tardive du champ lui-même.}
    \label{fig:turbu_avec_bulle}
\end{figure}

\begin{table}[H]
    \centering
    \caption{Grandeurs turbulentes sans et avec la bulle (moyennes stationnaires,
    niveau 8 dans les deux cas : champ ré-équilibré $t \ge 67$ vs run \texttt{lvl8}
    $t \ge 90$, sur la fenêtre complète).}
    \label{tab:turbu_avec_bulle}
    \begin{tabular}{@{}lccc@{}}
        \toprule
        \textbf{Grandeur} & \textbf{sans bulle} & \textbf{avec bulle} & \textbf{écart} \\
        \midrule
        Énergie cinétique $k$        & $15{,}0$  & $13{,}8$  & $-8\,\%$ \\
        Dissipation $\varepsilon$    & $1{,}36$  & $1{,}18$  & $-13\,\%$ \\
        $Re_\lambda$                 & $24{,}6$  & $24{,}4$  & $-1\,\%$ \\
        Échelle de Kolmogorov $\eta$ & $1{,}45$  & $1{,}51$  & $+4\,\%$ \\
        \bottomrule
    \end{tabular}
\end{table}

Une seule bulle occupe $\sim 0{,}1\,\%$ du volume du domaine ($\frac{4}{3}\pi R_0^3
\approx 2145$ pour $L_0^3 = 1{,}73\times 10^6$) : sans surprise, elle ne modifie la
turbulence \emph{globale} que modérément ($\lesssim 10\,\%$,
Tableau~\ref{tab:turbu_avec_bulle}), sans aucun saut à l'injection. La légère baisse
de $k$ et de $\varepsilon$ est en partie imputable à la traînée et à la tension de
surface locales, et en partie à l'évolution \emph{propre} du champ forcé (qui dérive
lentement en fin de run -- la référence sans bulle, elle-même stoppée dès l'atteinte
de la cible, n'était pas parfaitement stationnaire). On reste donc dans un régime de
\textbf{couplage essentiellement à sens unique} : la caractérisation de la turbulence
(\S\ref{sec:carac_we1}) reste \emph{représentative} des conditions vues par la bulle
pendant sa remontée, ce qui valide d'interpréter $U_b$ comme la réponse de la bulle
à une turbulence \emph{imposée}. (À l'échelle locale du sillage, un couplage à
double sens existe évidemment ; c'est seulement la statistique globale qui est peu
affectée.)
